% =====================================================================
% Paper one of the Lyrebird project: W0, benchmark inflation under
% recordist leakage.
%
% This is NOT the system paper. docs/24 §24.3 separates the two on
% purpose: the system, the resource and the 37 written refusals are
% ../arxiv/main.tex, and depend on the pipeline being correct. This
% paper depends on no part of that pipeline, which is what gives it the
% wider reach and is why it is a separate document.
%
% Build:  tectonic -X compile main.tex --outdir build
% =====================================================================
\documentclass[10pt,twocolumn,a4paper]{article}

\usepackage[T1]{fontenc}
\usepackage[utf8]{inputenc}
\usepackage{lmodern}
\usepackage[british]{babel}
\usepackage[margin=1.9cm]{geometry}
\setlength{\columnsep}{6mm}
\usepackage{amsmath,amssymb}
\usepackage{graphicx}
\usepackage{booktabs}
\usepackage{longtable}
\usepackage{array}
\usepackage{caption}
\usepackage{microtype}
\usepackage[numbers,sort&compress]{natbib}
\usepackage{xcolor}
\usepackage[colorlinks=true,allcolors=blue!55!black,breaklinks=true]{hyperref}

\captionsetup{font=small,labelfont=bf}
\captionsetup[table]{skip=4pt}

%: A visible marker for a number that is genuinely not yet in the repository.
\newcommand{\todo}[1]{{\color{red}\textbf{[TODO: #1]}}}
\newcolumntype{L}[1]{>{\raggedright\arraybackslash}p{#1}}

\title{\textbf{What a recordist-blocked split costs\\
a bioacoustic foundation model}}

\author{Mag.~art.~Shahab Nedaei\\
  \normalsize variable.gallery\\
  \normalsize \texttt{ned.tabulov@gmail.com}}

\date{4 August 2026}

\begin{document}
\makeatletter
\twocolumn[\begin{@twocolumnfalse}
\maketitle
\begin{abstract}
\noindent
Standard evaluation in bioacoustics splits recordings at random. Where a species'
recordings come disproportionately from one recordist, a random split puts recordings by
that person on both sides of the partition, and the reported accuracy is partly a measure of
recorder identity rather than of the bird. Holding recordists out is not a new idea:
\citet{stowell2014} did it twelve years ago and stated the rationale in the words we would
use. What has never been published is the \emph{quantity} --- what blocking costs a current
foundation model --- and without that number the practice stays easy to ignore.
%
Measuring it on an archive corpus is not straightforward, because both of the field's
archive-scale models were trained on the archives. Perch~2.0 trained on xeno-canto and
iNaturalist; BirdNET~V2.4 trained on xeno-canto; BirdSet excludes BirdNET from its own
evaluation for exactly this reason. A top-1 comparison on archive audio therefore partly
measures which recordings were memorised. We close that with two changes that also make the
experiment cheaper. The models are \textbf{frozen} and only a linear probe is trained, so
whatever either model memorised is identical on both sides of the comparison and cancels out
of the difference. And the audio was \textbf{uploaded after both training snapshots} ---
xeno-canto's \texttt{since:} tag filters on upload date, which is a one-directional
guarantee rather than an estimate. Because the discriminator is upload date rather than
archive provenance, ``reproduced on both models'' becomes possible on a xeno-canto corpus,
which the pre-registered design had ruled out.
%
On 8\,287 recordings uploaded after 2025-09-01, covering 884 species from 646 recordists, a
linear probe over frozen features loses \textbf{18.51} points of top-1 accuracy on Perch and
\textbf{21.59} on BirdNET when the split is blocked by recordist, and blocking is worse in 20 of
20 paired folds on both. As error rates that is 2.01$\times$ and 2.03$\times$ --- the two models
agree on the ratio to two decimal places while differing by three points in the raw figure. The
pre-registered bar was more than 10 points on both models, so the gate is met.
%
The drop is not a single number, and we say so: across nested label sets of 62 to 884 species on
the same recordings, the points climb from 12.75 to 18.51 while the error ratio falls from 2.97 to
2.01. Neither unit is invariant, so a figure quoted without its label-set size cannot be compared
against.
%
Two further measurements say what to do about it. The cost falls monotonically with how many
recordists a species has --- \textbf{30.8 points at two recordists against 3.1 at six or more},
with recordings per species and label-set size both held flat --- which turns ``more recordists
per species'' from a slogan into a target. And the \emph{same individual recordists} are the
findable ones in both models: per-recordist findability correlates at \textbf{+0.85} between
Perch and BirdNET after removing the effect of how many recordings each contributed, against a
shuffled null of +0.14.
%
Separately, and on identical rows, the two foundation models disagree about the size of the
recordist confound by more than either differs from a 32-dimensional space trained on one
small corpus: Perch sits 16.5\,\% below that space and BirdNET 24\,\% above it. If recorder
invariance were a product of archive-scale training, both would sit below. It is better read
as a property of the archives that any model trained on them inherits.
%
The deliverable is not a system. It is a number, a threat model and the code: block your
splits by recordist, here is what it costs, here is how to reproduce it.
\end{abstract}
\vspace{4mm}
\end{@twocolumnfalse}]
\makeatother

% =====================================================================
\section{Introduction}
\label{sec:intro}

Automatic species classification from sound is used for real monitoring and real
conservation decisions. BirdNET \citep{birdnet2021} and Perch \citep{perch2} are deployed by
ecologists, and their benchmark numbers set expectations about what a deployment will do.

Those numbers come from evaluations that split recordings at random. That is only sound if a
recording's provenance carries no information about its label. In the archives it does. A
species' recordings are contributed by volunteers, and for many species they come
disproportionately from a handful of people. A recordist is not only a person: they are a
microphone, a preamplifier, a set of recording habits, a geographic range and a preferred
time of day. When recordings by one person land on both sides of a random split, a model can
reach the right species by recognising the person.

This matters beyond the leaderboard, because the failure has a deployment shape. A recordist
is functionally \emph{a recorder and a site}. Generalising across recordists is the same
problem as generalising to a new device in a new place, which is the transition the field
already knows disappoints \citep{perch2}. Naming a measurable cause, and pricing it, is more
actionable than the observation that deployment is hard.

\paragraph{What this paper contributes.} One number, measured carefully, plus the method that
makes it measurable at all. Specifically:

\begin{enumerate}
  \item \textbf{The cost of blocking}, for a linear probe over frozen features from two
    independent foundation models, on audio neither model can have trained on
    (\S\ref{sec:results}).
  \item \textbf{A contamination-free design} for asking the question on archive audio, which
    the obvious design cannot do (\S\ref{sec:design}). This is what lets both models be
    tested on the same corpus.
  \item \textbf{The observation that the pre-registered gate was written in the wrong unit.}
    Accuracy points are not comparable across label-set sizes, and we show the drop's
    dependence on class count directly (\S\ref{sec:sweep}).
  \item \textbf{A dose-response curve, and therefore a target} (\S\ref{sec:dose}): the cost falls
    from 30.8 points at two recordists per species to 3.1 at six or more, with label-set size and
    recordings per species held flat. This is the only result here that removes the problem
    instead of pricing it.
  \item \textbf{Evidence that the same individual recordists are findable in both models}
    (\S\ref{sec:samepeople}): per-recordist findability correlates at $+0.85$ between Perch and
    BirdNET after the number of recordings is partialled out, against a shuffled null of $+0.14$.
    The structure belongs to those recordists, not to either architecture.
  \item \textbf{Evidence that the structure is in the archives, not in one model}
    (\S\ref{sec:archives}): read on identical rows, Perch and BirdNET disagree about the size
    of the recordist confound by more than either differs from a 32-d space trained on one
    small corpus. Recorder invariance is not conferred by archive-scale training.
  \item \textbf{Code and a fixed population}, so the number can be checked rather than
    believed (\S\ref{sec:repro}).
\end{enumerate}

\paragraph{What it does not contribute.} Not the practice of recordist-blocked evaluation,
which is twelve years old. Not a new model, not a new architecture, and not a claim that
these models are bad. A model that reads recorder identity is doing something reasonable
with the data it was given; the problem is the evaluation protocol that rewards it.

% =====================================================================
\section{The practice is old, the quantity is missing}
\label{sec:related}

\paragraph{The practice exists, and is older than the models.} \citet{stowell2014} split by
recordist and gave the rationale plainly: the approach ``tests whether the system can
generalise to recordings from unknown recordists, rather than adapting to any specifics of
the known recordists''. Read against the full text rather than the abstract, two details
matter here. The stratification was applied to \emph{one} of their four datasets ---
\texttt{bldawn}, 20 recordings each from three UK recordists, three-fold, training on two
recordists and testing on the held-out third. And the paper reports \textbf{no
random-against-blocked comparison}: recordist stratification is used where it is used, and
its cost is never priced. The Cornell CBI dataset is likewise partitioned so that no
recordist spans two partitions.

\paragraph{The quantity does not exist.} We checked this against sources rather than
abstracts. Neither the 2026 foundation-model comparative review
(\href{https://arxiv.org/abs/2508.01277}{arXiv:2508.01277}) nor the 2025 large-scale training
recipe (\href{https://arxiv.org/abs/2508.11845}{arXiv:2508.11845}) reports a
random-against-blocked accuracy comparison. Nobody has published what a recordist-blocked
split costs a current foundation model.

\paragraph{So the claim moves.} The deliverable is not \emph{block your splits}, which is
known and ignored. It is \emph{here is what one costs}. That is a smaller claim than the
project originally registered, and it is the one that is actually unoccupied. It is also the
claim most likely to change behaviour, because the reason a twelve-year-old
recommendation gets ignored is that nobody has priced ignoring it.

% =====================================================================
\section{Why the obvious experiment is invalid}
\label{sec:design}

The experiment as first designed was: take a public archive-derived corpus, take Perch and
BirdNET off the shelf, evaluate top-1 species accuracy twice, changing only the split.
That design does not work, and the reason is worth stating because it applies to any
evaluation of these models on archive audio.

\paragraph{Both models trained on the archives.} Perch~2.0's technical report
\citep{perch2} lists roughly 1.5\,M recordings over 14\,500 classes drawn from xeno-canto,
iNaturalist research-grade audio via GBIF, the Tierstimmenarchiv and FSD50K. BirdNET~V2.4
\citep{birdnet2021} trained on xeno-canto. BirdSet already excludes BirdNET from evaluation
on its \texttt{XCM}/\texttt{XCL} splits on precisely these grounds, since those splits are
drawn from xeno-canto.

\paragraph{Why that is fatal rather than merely untidy.} A memorised recording is scored
correctly for reasons that have nothing to do with acoustics, and a model that memorised a
recording will also place it near that recording's siblings. So on archive audio, top-1
accuracy measures memorisation as much as generalisation --- and the \emph{gap} between two
splits partly measures which recordings were memorised, rather than what blocking costs.
Under the pre-registered design this made ``reproduced on both models'' unreachable on any
xeno-canto-derived corpus.

\paragraph{Two changes fix it, and both make the experiment cheaper.}

\begin{enumerate}
  \item \textbf{Freeze the model; train only a linear probe.} The foundation model is a fixed
    feature extractor in both arms. Whatever it memorised is therefore present identically on
    both sides of the comparison and cancels out of the difference. What differs between the
    arms is only the probe's training set, which is the thing the claim is about.
  \item \textbf{Use audio uploaded after both training snapshots.} Xeno-canto's
    \texttt{since:} tag filters on \emph{upload} date. A recording uploaded after a model's
    data snapshot cannot have been in that snapshot regardless of when it was recorded, which
    makes it a one-directional guarantee rather than an estimate. We set the bound to
    2025-09-01: one month past Perch~2.0's publication, and more than two years past
    BirdNET~V2.4's June~2023 release. No attempt at precision is made, because the argument
    needs the bound to be safe rather than tight.
\end{enumerate}

\paragraph{This is what makes the two-model requirement satisfiable.} Because the
discriminator is upload date and not archive provenance, the objection that disqualifies
BirdNET on a xeno-canto corpus does not apply: the audio postdates its training set by two
years. The design constraint that broke the registered method is dissolved rather than worked
around, and both models can be read on identical rows.

\paragraph{What the design still cannot claim.} A linear probe over frozen features is not a
fine-tuned model, and the absolute accuracies here are not comparable to a published
benchmark number. That is deliberate: the quantity of interest is a \emph{difference} between
two splits under an otherwise identical procedure, and the frozen-feature design is what makes
that difference attributable to the split. Whether full fine-tuning would widen or narrow the
gap is open, and \S\ref{sec:threats} says so.

% =====================================================================
\section{Data}
\label{sec:data}

\paragraph{The query, and a warning that belongs in it.} The population is every xeno-canto
recording matching \texttt{grp:birds since:2025-09-01 q:A}. The group tag is not optional.
Xeno-canto hosts bats, frogs, grasshoppers and land mammals alongside birds, and the same
query without \texttt{grp:birds} returned 6\,000 recordings of which \textbf{none were
birds} --- 3\,035 bats, 1\,568 frogs, 926 grasshoppers and 471 land mammals. Pagination
simply never reached the birds. Nothing errored, and the first downloads were 256\,kHz
ultrasonic bat clips that decoded perfectly well and would have entered a bird model as
though they were birdsong. Any replication must carry the group tag.

\paragraph{Census.} The complete query returns \textbf{21\,928} recordings, of which
21\,709 have downloadable audio; xeno-canto redacts the file for species it considers at
risk from trapping or disturbance, and no query gets round it. Those recordings cover
\textbf{4\,305} species and \textbf{961} recordists.

\paragraph{The species floor is the experiment, not a convenience.} Under a
recordist-blocked split, a species whose recordings all come from one recordist lands
entirely on one side of the partition and cannot be classified at all. That is an empty test
class, not leakage, and scoring it as an error would manufacture the result the experiment is
looking for. We therefore keep only species with at least 6 recordings from at least
\textbf{2 distinct recordists}, which is \textbf{886} species covering 13\,646 recordings,
and we report the number excluded.

\paragraph{A cap of 12 recordings per species, which improves the experiment.} The scorable
population is severely skewed: 6 recordings at the floor, a median of 9, and a maximum of
338. Left uncapped, a handful of species would supply a large share of the probe's training
rows and the class prior would do work that the features are supposed to do. Capping at 12
reduces the imbalance from 56:1 against the floor to 2:1, loses \emph{no} species because the
floor is 6, and reduces the download by roughly 5\,300 recordings.

The cap has to be recordist-aware. Consider a species with 20 recordings where one recordist
contributed 19 and a second contributed 1: an ordinary random sample of 12 drops that second
recordist about 40\,\% of the time, the species falls to a single recordist, and it becomes
unscorable under a blocked split \emph{after} having been counted as scorable. Selection
therefore round-robins across a species' recordists, which keeps every recordist for as long
as the cap allows, and an assertion checks that the scorable species set is unchanged by
capping.

\paragraph{Selected, then scored.} The cap selects \textbf{8\,321} recordings over 886 species
and 649 recordists, with 6 to 12 recordings per species and at least 2 recordists each. Of those,
24 were lost --- 13 that no decoder would read and 11 that returned an HTTP error --- so the
population \emph{scored} is \textbf{8\,287} recordings over \textbf{884} species from
\textbf{646} recordists. No species fell below the floor as a result of those losses.

We keep the two figures apart because the pilot did not. It fetched 2\,158 recordings and is
usually described by that number, but only \textbf{730} cleared its floor and were scored, over
62 species on one model. Table~\ref{tab:pilot} sets the two runs side by side.

% =====================================================================
\section{Method}
\label{sec:method}

\paragraph{Features.} Each recording is decoded, mixed to mono, resampled, and cut into the
model's native window: 5\,s at 32\,kHz for Perch (six windows, 30\,s) and 3\,s at 48\,kHz
for BirdNET (ten windows, the same 30\,s). Window embeddings are mean-pooled to one vector
per recording --- 1536-d for Perch, 1024-d for BirdNET's backbone. Both models are ONNX
exports run without gradients; neither is updated at any point.

\paragraph{Probe.} A multinomial logistic regression with an L2 penalty
($C = 1$), on features standardised using the \emph{training} side's mean and variance only.
Standardising over both sides would leak the test set's statistics across the partition,
which is the same class of error the paper is about.

\paragraph{The probe runs on the GPU, and that required a check.} At 884 classes in 1536
dimensions a full-batch multinomial fit costs on the order of $10^{10}$ multiply--adds per
iteration, which is not tractable on CPU across 20 folds, two arms and two models. The probe
is therefore implemented in PyTorch with full-batch L-BFGS, minimising the same objective
scikit-learn minimises. Swapping the estimator silently would make this run incomparable with
the pilot, so the two are checked against each other on synthetic data and on the pilot's real
features; they agree to within one accuracy point, and the check fails loudly if they do not.

\paragraph{Three arms.} All three restrict each fold to the classes present on
\emph{both} sides before scoring, and that restriction must be applied to every arm or the
comparison is rigged: a blocked split strands whole species on one side, and dropping those
from the blocked arm alone would leave it a smaller and easier label set than the random arm,
which inflates blocked accuracy and understates the very drop being measured.

\begin{itemize}
  \item \textbf{A --- random.} 30\,\% of each species' recordings held out, stratified by
    species. This is the standard protocol.
  \item \textbf{B --- recordist-blocked.} 30\,\% of recordists held out; every recording by a
    held-out recordist is on the test side.
  \item \textbf{C --- paired, and this is the number we quote.} Each blocked fold's surviving
    class set is reused for a random fold at the same held-out fraction, so the two differ in
    exactly one thing: whether a recordist may appear on both sides. Arms A and B are not
    scored on the same task, because blocked folds end up with fewer classes, and fewer
    classes is an easier problem.
\end{itemize}

Twenty folds per arm, seeded. Because the folds are paired, the per-fold difference is the
statistic with the least noise in it, and we report its spread and the number of folds in
which blocking was worse, not only the mean.

\paragraph{Reporting the effect in two units.} The pre-registered gate is a drop of more than
10 points of top-1 accuracy, reproduced on both models, and it is scored as written. But
points are a poor unit near the ceiling: the pilot's 9.68 points at 93\,\% accuracy was
$2.44\times$ the errors, whereas the same 9.68 points at 60\,\% accuracy would be $1.24\times$.
We therefore report the error-rate ratio beside the gate, and \S\ref{sec:sweep} shows what the
points figure does as the label set grows.

% =====================================================================
\section{Results}
\label{sec:results}

The population scored is \textbf{8\,287} recordings over \textbf{884} species from
\textbf{646} recordists. Of the 8\,321 selected, 24 were lost --- 13 that no decoder would
read and 11 that returned an HTTP error --- and no species fell below the floor as a result.

\paragraph{The gate is met, on both models, and not narrowly.} A recordist-blocked split costs
a linear probe over frozen Perch~2.0 features \textbf{18.51 points} of top-1 accuracy, and one
over frozen BirdNET~V2.4 features \textbf{21.59 points}. The pre-registered bar was more than
10 points reproduced on both models. Table~\ref{tab:headline} is the paired arm, which is the
one to quote.

\begin{table}[t]
\footnotesize
\setlength{\tabcolsep}{4pt}
\caption{The paired arm: each fold's blocked class set reused for a random fold, so the two
differ only in whether a recordist may appear on both sides. 20 folds, 664 classes a fold.}
\label{tab:headline}
\begin{tabular}{L{2.5cm} r r}
\toprule
 & Perch 2.0 & BirdNET V2.4 \\
 & 1536-d & 1024-d \\
\midrule
random split      & 81.68\,\%  & 79.08\,\% \\
recordist-blocked & 63.16\,\%  & 57.49\,\% \\
\textbf{drop}     & \textbf{18.51} & \textbf{21.59} \\
per-fold sd       & 2.78       & 2.96 \\
blocked worse in  & 20/20      & 20/20 \\
error rate        & 18.32 $\to$ 36.84\,\% & 20.92 $\to$ 42.51\,\% \\
as a ratio        & 2.01$\times$ & 2.03$\times$ \\
\midrule
gate $>10$ points & \textbf{met} & \textbf{met} \\
\bottomrule
\end{tabular}
\end{table}

Three things about the shape of that result. Blocking was worse in \textbf{20 of 20 paired
folds} for both models, with a per-fold standard deviation under 3 points, so this is not a
mean hiding a mixed picture. The two models agree on the error ratio to two decimal places
(2.01 and 2.03) while differing by 3 points in the raw figure, which is the sense in which the
ratio is the more transferable statistic. And BirdNET starts lower and falls further, so the
ordering of the two models is the same under both splits --- blocking does not reverse which
model looks better.

\paragraph{The pilot understated it, and the reason is its population.} The pilot measured 9.68
points on 62 species and recorded that as failing the gate. At 62 classes drawn from this
population the same procedure gives \textbf{12.75} points on Perch. The remainder of the gap
between 9.68 and 18.51 is the label-set size, which \S\ref{sec:sweep} takes apart.

\subsection{The drop depends on how many species are in the label set}
\label{sec:sweep}

Table~\ref{tab:sweep} runs the paired arm on nested label sets drawn from the same recordings
with the same folds, so the only thing changing is how many species the probe must separate.

\begin{table}[t]
\footnotesize
\setlength{\tabcolsep}{5pt}
\caption{The same experiment at five label-set sizes. Points climb with class count. The error
ratio falls. \textbf{Neither is invariant}, which is a caution against quoting either alone.}
\label{tab:sweep}
\begin{tabular}{r rr rr}
\toprule
 & \multicolumn{2}{c}{Perch} & \multicolumn{2}{c}{BirdNET} \\
\cmidrule(lr){2-3}\cmidrule(lr){4-5}
species & points & ratio & points & ratio \\
\midrule
62  & 12.75 & 2.97$\times$ & 16.07 & 2.93$\times$ \\
150 & 15.28 & 2.66$\times$ & 19.05 & 2.65$\times$ \\
300 & 17.15 & 2.40$\times$ & 20.61 & 2.39$\times$ \\
500 & 18.76 & 2.23$\times$ & 21.62 & 2.23$\times$ \\
884 & 18.51 & 2.01$\times$ & 21.59 & 2.03$\times$ \\
\bottomrule
\end{tabular}
\end{table}

We predicted that the ratio would be flat and the points would climb. \textbf{Half of that is
wrong.} The points do climb, steeply from 62 to 500 and then flat. But the ratio falls
monotonically too, from about 2.95 to about 2.02. So there is no unit in which this effect is
a single number independent of the label set, and a paper reporting one figure without its
class count is not saying enough to be compared against.

What survives is narrower and still useful. The two models track each other almost exactly at
every size --- the ratio columns agree to within 0.02 at four of the five sizes --- so the
label-set dependence is a property of the task and not of a model. And a bar written in
accuracy points is the more treacherous of the two: our own pre-registered 10-point gate is
failed at 62 species by the pilot and cleared at every size here, on the same effect.

\paragraph{So the gate was a bad instrument that happened to give the right answer.} It should
have been registered as an error-rate ratio with the label-set size stated. We are not moving
it after the fact, and we are recording that it was written in the wrong unit.

% =====================================================================
\section{How many recordists per species does it take?}
\label{sec:dose}

The recommendation that follows from a leakage result is normally ``block your splits''. That
prices the problem without fixing it. If the leakage comes from a species being carried by too
few people, the fix is more people --- and that is only actionable with a number attached.

Table~\ref{tab:dose} measures the paired drop separately for species carried by 2, 3, 4, 5 and
6 or more recordists. Two confounds are held down. Every stratum is subsampled to the same
\textbf{115 species}, so a drop in points is comparable between them. And the per-species cap of
12 keeps recordings per species between 8.2 and 10.9 across all five strata, so the trend is not
simply more training rows.

\begin{table}[t]
\footnotesize
\setlength{\tabcolsep}{5pt}
\caption{The cost of blocking, by how many recordists carry the species. 115 species a stratum,
recordings per species held flat. \textbf{The effect falls by a factor of ten.}}
\label{tab:dose}
\begin{tabular}{c rr r}
\toprule
recordists & drop & error ratio & rec./species \\
\midrule
2        & 30.82 & 5.83$\times$ & 8.2 \\
3        & 20.92 & 3.16$\times$ & 8.6 \\
4        & 18.52 & 2.59$\times$ & 8.2 \\
5        & 11.12 & 2.05$\times$ & 8.9 \\
\textbf{6+} & \textbf{3.08} & \textbf{1.29$\times$} & 10.9 \\
\bottomrule
\end{tabular}
\end{table}

\paragraph{At six recordists the leakage is essentially gone.} 3.08 points and 1.29 times the
errors, against 30.82 points and 5.83 times at two. The decline is monotonic and steepest between
4 and 6. The 6+ stratum does carry slightly more recordings per species (10.9 against about 8.4),
so a small part of its advantage may be sample size --- but the fall from 30.82 to 11.12 happens
across strata whose recordings per species are 8.2, 8.6, 8.2 and 8.9, which is flat. Sample size
is not what moves this.

\paragraph{What that licenses as advice.} A species already covered by six or more independent
recordists is close to safe under a random split. A species covered by two is not usable for
benchmarking without blocking. And for an archive, one additional recording of a two-recordist
species by a \emph{third} person is worth far more than a fourth recording by one of the existing
two. That is a concrete collection priority, and it is the only intervention here that removes the
problem instead of measuring it.

% =====================================================================
\section{The same recordists are findable in both models}
\label{sec:samepeople}

Both models carrying an aggregate recordist lift is weak evidence for a mechanism, because two
models can show the same aggregate for unrelated reasons. The sharper question is whether the
\emph{same individuals} are the easy ones.

For each of the 212 recordists with at least 6 recordings in the population, we measure
findability: how much more often a recording's nearest neighbour by cosine is also theirs than
chance allows. Then we correlate that per-recordist score between the two models.

\begin{table}[t]
\footnotesize
\setlength{\tabcolsep}{5pt}
\caption{Per-recordist findability, Perch against BirdNET, over 212 recordists. The partial
correlation is the headline: it removes the number of recordings each contributed, which inflates
findability in both spaces at once.}
\label{tab:samepeople}
\begin{tabular}{L{4.6cm} r}
\toprule
mean findability, Perch & $+0.361$ \\
mean findability, BirdNET & $+0.438$ \\
\midrule
Pearson, raw & $+0.874$ \\
Spearman, raw & $+0.880$ \\
\textbf{Pearson, log(recordings) removed} & $\mathbf{+0.852}$ \\
\midrule
shuffled null, 95th percentile & $+0.140$ \\
findability vs.\ log(recordings) & $+0.37$, $+0.39$ \\
positive in both spaces & 83.0\,\% \\
\bottomrule
\end{tabular}
\end{table}

\paragraph{The nuisance variable, and why the partial correlation is the number.} A recordist with
20 recordings is easier to find than one with 6 in \emph{either} space, because more same-label
neighbours exist to be found. That inflates both scores together and would manufacture agreement
with no shared mechanism behind it. The effect is real --- findability correlates with
log(recordings) at about $+0.38$ in both spaces --- and it accounts for almost none of the
agreement. Partialling it out moves the correlation from $+0.874$ to $+0.852$, still six times the
shuffled null's 95th percentile.

\paragraph{What this settles.} Two architectures, trained by different groups on overlapping
archives, agree at $+0.85$ about \emph{which people} are identifiable, and 83\,\% of recordists are
identifiable in both. The structure is a property of those recordists --- their microphone, their
sites, their habits --- and not of either model's inductive bias. A model cannot be blamed for it
and, as \S\ref{sec:archives} shows, cannot be trained out of it either. \S\ref{sec:dose} is where
the remedy is.

% =====================================================================
\section{The structure is in the archives, not in one model}
\label{sec:archives}

A drop measured on one model could be a property of that model. The measurements in this
section are a separate experiment, on a separate population, and they are what rules that
out. They use nearest-neighbour retrieval lift against a permutation null rather than a
probe: for a given grouping, how much more often is a recording's nearest neighbour in its own
group than chance allows.

Three embedding spaces are compared on \emph{identical rows} with an identical estimator ---
Perch~2.0 at 1536-d, BirdNET~V2.4 at 1024-d, and a 32-d task-specific space trained on this
project's own corpus (the ``critic'', described in the companion system paper). Two
populations: one drawn to be recordist-stratified, one drawn as species pairs.

\begin{table}[t]
\footnotesize
\setlength{\tabcolsep}{4pt}
\caption{Nearest-neighbour lift over a permutation null, same rows, same estimator. Every
space is confounded in the same direction; they differ only in degree. \textbf{The two
foundation models disagree with each other by more than either differs from a 32-d space
trained on one small corpus}, which is the observation this section exists for.}
\label{tab:crossmodel}
\begin{tabular}{L{1.5cm} L{1.5cm} r r}
\toprule
population & space & species & recordist \\
\midrule
T1        & critic 32-d   & ---      & $+0.0891$ \\
1\,567 rec.\ & Perch 1536-d  & ---      & $+0.0744$ \\
135 rec'ists & BirdNET 1024-d & ---     & $+0.1107$ \\
\midrule
T3        & critic 32-d   & $+0.0074$ & $+0.1684$ \\
1\,644 rec.\ & Perch 1536-d  & $+0.0588$ & $+0.1155$ \\
760 rec'ists & BirdNET 1024-d & $+0.0393$ & $+0.1755$ \\
\bottomrule
\end{tabular}
\end{table}

\paragraph{Every space carries it, and the two big models disagree about how much.} On the
recordist-stratified population, Perch sits \textbf{16.5\,\% below} the small
corpus-specific space and BirdNET \textbf{24\,\% above} it. On the species-pair population
Perch is 31.4\,\% below and BirdNET 4\,\% above. Two models trained by different groups on
overlapping archives, read on the same recordings by the same estimator, differ from each
other more than either differs from a 32-d space trained on one project's corpus.

\paragraph{What that licenses, and what it does not.} It licenses the claim that
\textbf{recorder invariance is not conferred by archive-scale training}. If it were, both
archive-scale models would sit below the small space, and one of them sits above it. The
structure is better explained as a property of the archives that any model trained on them
inherits to some degree. It does \emph{not} license a ranking of the two models: the
populations here are small, and the T3 recordist column rests on a coverage of only 0.329
because most recordists in that draw have too few recordings for the estimator's floor.

\paragraph{And the structure survives on audio the model never saw.} Repeating the
recordist measurement on the post-cutoff population of \S\ref{sec:data} --- rows Perch
provably did not train on --- gives a recordist lift of $+0.441$ against a species lift of
$+0.750$ on 293 matched rows where both groupings clear the estimator's floor, both at the
permutation floor for significance. Whatever the recordist structure is, it is not Perch
recognising its own training data.

\paragraph{A caution about a number we have retracted elsewhere.} An earlier and louder
version of this project's headline was a recordist lift of $+0.488$ against a species lift of
$+0.196$, at \emph{syllable} level. A benchmark split is a \emph{recording}-level operation,
and at recording level on the same corpus the ordering reverses: species $+0.229$ against
recordist $+0.196$. That reversal is itself a selection artefact --- the estimator's
eight-member floor at recording level means ``species with at least eight recordings'', which
is 77 of 3\,636, and recordings per species correlates with recordists per species at
$r = +0.923$. The honest reading is that the recording-level figures are the ones relevant to
a split, that they are smaller than the syllable-level ones, and that the population which
makes them measurable is the population where the confound is weakest. The pre-registered
10-point gate should not have been treated as likely on the strength of $+0.488$.

% =====================================================================
\section{Recordist is not the only axis, and it is not the most expensive}
\label{sec:axes}

Everything above prices one blocking variable. A benchmark designer chooses among several, and
two others are standard: \textbf{geography}, because train-here-deploy-there is the deployment
case, and \textbf{time}, because equipment and archive conventions drift. Both are answerable on
this population without re-embedding anything, since every recording carries a country, a
date, and --- where xeno-canto publishes it --- a latitude and longitude.

Everything is held identical to \S\ref{sec:results}: the same 8\,287 recordings, the same 884
species, the same frozen embeddings, the same probe, the same held fraction, the same seeds and
the same paired class-set correction. Only the blocking variable changes. The recordist arm is
therefore a replication check on this harness, and it reproduces: $18.38$ against the $18.51$
of \S\ref{sec:results}.

\paragraph{The unit has to be recordings, not levels.} Holding out a fixed share of blocking
\emph{levels} is right when comparing one variable against a random split, and wrong when
comparing variables against each other, because the variables differ enormously in how many
levels they have. Thirty per cent of 116 countries strands 43\,\% of the recordings; thirty per
cent of 646 recordists strands 34\,\%. The arm with less training data then loses accuracy for
a reason that has nothing to do with leakage. So each arm below adds shuffled levels until the
held \emph{recordings} reach a common target --- the same correction in spirit as the class-set
pairing that \S\ref{sec:method} applies.

\begin{table}[t]
\footnotesize
\setlength{\tabcolsep}{4pt}
\caption{Top-1 drop by blocking variable, held recording share matched across arms, 20 paired
folds. Every axis costs far more than the pre-registered 10-point gate, and country costs the
most on both models.}
\label{tab:axes}
\begin{tabular}{L{2.5cm} r r r r}
\toprule
blocked by & levels & classes & Perch & BirdNET \\
\midrule
recordist & 646 & 716 & $19.02$ & $22.07$ \\
location, $1^\circ$ cell & 1\,128 & 726 & $18.26$ & $20.68$ \\
\textbf{country} & 116 & 502 & $\mathbf{23.87}$ & $\mathbf{26.97}$ \\
year & 36 & 706 & $19.11$ & $21.50$ \\
\bottomrule
\end{tabular}
\end{table}

\paragraph{These axes are not independent, and that is the point.} A recordist records near
home, so a location-blocked split is substantially a recordist-blocked split. Measured as
normalised mutual information, recordist against $1^\circ$ cell is $0.817$ and recordist
against country is $0.733$; recordist against year is only $0.246$. So the transferable claim
is \emph{not} that there are four separate leaks to close. It is that blocking costs 17--27
points on whichever of these axes a benchmark chooses, and that a geographically blocked
benchmark is the most expensive of the four.

\paragraph{Why the country result is reported as suggestive.} Two biases oppose each other in
it and neither can be removed on this population. Even matched, country holds $0.436$ of the
recordings against recordist's $0.402$, because 116 coarse levels cannot be matched finely, and
that residual favours a larger country drop. Against it, country's folds survive with far fewer
classes (502 against 716), and \S\ref{sec:sweep} established that a smaller label set makes the
drop \emph{smaller}. Sub-sampling within countries to match the held share exactly, or
re-running every arm at a common class count, would settle the ordering. We report the numbers
with both biases named rather than the ordering as a finding.

\paragraph{What does survive without qualification} is the weaker and more useful statement:
there is no cheap axis. A designer hoping that blocking by place or by year avoids the cost of
blocking by recordist does not get that discount --- year costs what recordist costs, and place
costs at least as much.

% =====================================================================
\section{Threats to validity}
\label{sec:threats}

\paragraph{One archive, one quality tier.} Everything here is xeno-canto at \texttt{q:A}.
Quality~A keeps the audio comparable but selects recordings that a human judged clean, and
cleanliness may correlate with recordist skill --- which is the variable under test. A
replication across quality tiers would say whether the effect is larger or smaller in the
messier part of the archive.

\paragraph{Upload date is not recording date.} The guarantee is one-directional and that is
all it needs to be: a recording uploaded after a snapshot was not in that snapshot. It does
not follow that these recordings are acoustically representative of the archive as a whole,
and a corpus of recent uploads may over-represent currently active recordists.

\paragraph{A probe is not a fine-tuned model.} The frozen-feature design is what makes the
difference attributable to the split, at the cost of not being the number a practitioner
would report from fine-tuning. The direction of that difference is not predictable from this
experiment.

\paragraph{Recordist identity is self-reported metadata.} One person with two accounts reads
as two recordists and weakens blocking; two people sharing an account reads as one and
strengthens it. Both push in known directions but neither is quantified here.

\paragraph{Neither unit is invariant, including ours.} \S\ref{sec:sweep} shows the points
climbing and the ratio falling as the label set grows, so our own headline is a statement about an
884-class problem and nothing more general. We predicted the ratio would be flat and it is not.
A reader comparing against a benchmark of a different size should treat both figures as
size-dependent and quote the class count alongside either.

\paragraph{The dose curve is observational.} Species are not assigned recordists at random. A
species carried by six recordists may differ from a two-recordist species in ways that also affect
difficulty --- how common it is, how loud, how widely distributed. Holding label-set size and
recordings per species flat removes the two most obvious of those, not all of them. The curve is
strong enough to act on and not strong enough to call causal.

% =====================================================================
\section{Recommendations}
\label{sec:what}

\paragraph{For benchmark maintainers.} Report both numbers. A random split and a
recordist-blocked split, with the recordist coverage of the corpus stated, so that readers can
see the gap rather than infer it. Blocking need not replace random splitting to be useful;
the gap is the informative quantity.

\paragraph{For practitioners.} Where a deployment will meet new recorders in new places ---
which is most deployments --- the blocked figure is the closer estimate, and the gap measured
here is the scale of the optimism in the random one.

\paragraph{For corpus builders, and this is the actionable one.} Count recordists per species,
not recordings per species. \S\ref{sec:dose} puts the threshold at about six: below it a species
carries leakage worth 11 to 31 points, at or above it about 3. A species covered by two people is
not usable for benchmarking without blocking, however many recordings it has. When choosing what
to collect, one recording of a thinly-covered species by a \emph{new} recordist is worth more than
several more by an existing one.

\paragraph{And keep the metadata.} Recordist identity is what makes any of this measurable.
Archives that keep it, and evaluations that use it, cost little and are the precondition for
pricing this failure at all.

% =====================================================================
\section{Reproducibility}
\label{sec:repro}

Every number is produced by a script that prints its pre-registered gate before its result and
writes a JSON sidecar beside itself. The catalogue is cached, so the population is fixed
rather than re-searched.

\begin{table}[t]
\footnotesize
\setlength{\tabcolsep}{4pt}
\caption{The pilot and this run. The pilot's limits were not properties of the archive: its
catalogue held exactly 6\,000 recordings because its search stopped at 60 pages, and it
selected by \emph{recordist} because it inherited that rule from a different experiment whose
unit was the recordist.}
\label{tab:pilot}
\begin{tabular}{L{2.6cm} L{2.2cm} L{2.4cm}}
\toprule
 & pilot & this run \\
\midrule
catalogue read      & 6\,000        & 21\,928 \\
selection unit      & recordist     & species \\
recordings embedded & 2\,158        & 8\,321 \\
recordings scored   & 730           & 8\,287 \\
species scored      & 62            & 884 \\
recordists          & 137           & 646 \\
models              & Perch only    & Perch + BirdNET \\
probe               & CPU, sklearn  & GPU, verified equivalent \\
\bottomrule
\end{tabular}
\end{table}

\begin{table}[t]
\footnotesize
\setlength{\tabcolsep}{4pt}
\caption{What reproduces what. Filenames are relative to \texttt{pipeline/analysis/} in the
project repository; the \texttt{analysis/} prefix is dropped here because it does not fit a
two-column measure.}
\label{tab:repro}
\begin{tabular}{L{3.6cm} L{3.9cm}}
\toprule
script & what it does \\
\midrule
\texttt{w0\_full\_catalogue.py}
  & pages the complete post-cutoff catalogue and reports what it can support \\
\texttt{w0\_full\_fetch.py}
  & selects the scorable population, downloads it, embeds it through both models,
    checkpointed and resumable \\
\texttt{w0\_probe.py}
  & the GPU probe, and \texttt{-{}-verify} for its equivalence to scikit-learn \\
\texttt{w0\_full.py}
  & the three arms, the gate, and the class-count sweep \\
\texttt{w0\_recordist\_split.py}
  & the earlier pilot, kept runnable and unchanged \\
\bottomrule
\end{tabular}
\end{table}

An \texttt{XC\_API\_KEY} is required for the fetch; xeno-canto API v3 needs a personal key.
The embeddings are cached, so the analysis can be re-run without re-downloading, and
\texttt{w0\_recordist\_split.py -{}-self-test} checks the split machinery against two
synthetic worlds with known answers --- one where recordist is assigned at random, where
blocking must cost nothing, and one where recordist is confounded with species, where
blocking must be devastating.

% =====================================================================
\section*{Data and code availability}

\todo{Zenodo DOI for the archived release.} No audio is redistributed: the fetch scripts
retrieve it from xeno-canto under the contributors' own licences, and the cached catalogue
records exactly which recordings were used.

\section*{Acknowledgements}

The recordings this paper is about were contributed to xeno-canto \citep{xenocanto} by
volunteers, and the recordist metadata that makes the measurement possible is their work. It
is worth being precise about the debt: this paper reports a confound carried by *who* recorded
each file, and that confound is only measurable because those contributors are named. An
archive that discarded attribution would have hidden the effect rather than prevented it.

The measurement also depends on Perch and BirdNET being released with usable weights. A
claim that an evaluation protocol is optimistic is worth much less when the models under test
cannot be run by the people being asked to change the protocol, and both can.

\section*{Funding}

None. This is independent work by an artist, made at the boundary of art and science, with no
grant, institutional or commercial support of any kind.

\bibliographystyle{plainnat}
% `references`, not `../references/references`. A preprint server compiles from the root of
% the submission directory, so a path that climbs above it resolves to nothing and the
% bibliography comes out empty. An identical copy sits next to this file, so this form builds
% locally and on a submission server alike.
\bibliography{references}

\end{document}
